\chapter{结论与展望}
\label{cha:conclusion}

\section{全文总结}

本文围绕辅助交互与便携式使用场景下的低成本无创脑机接口系统展开研究。其出发点在于：面向这类实际使用需求，系统往往需要同时兼顾穿戴便利性、无线通信、持续运行与本地响应能力，因此难以直接沿用实验室中的高密度有线配置。由此带来的少导联采集、有限算力和在线交互约束，使得运动想象脑电解码任务需要从标准离线设置重新收束到更可部署的系统方案。围绕这一问题，本文从硬件系统、解码方法和在线联动实现三个层面开展了工作，并形成了如下结论。

第一，本文完成了面向双导便携式场景的脑电采集与基础交互系统设计与实现。以 KS1092 双通道前端和 nRF52832 主控为核心，构建了脑电采样、参数配置、BLE 回传和上位机显示的完整基础链路。通过 Alpha 节律验证和标签链路响应实验，可以确认该平台已经具备采集具有神经生理意义信号以及支持在线范式事件标记的能力，为后续在线试跑提供了系统基础。

第二，本文提出并完善了面向低通道端侧场景的改进型自注意力脑电解码方法 EdgeMIFormer。在标准 22 导条件下，模型完成了结构设计、性能验证以及剪枝量化后的部署验证；在低通道扩展验证中，统一协议结果表明，随着导联数由 22 导逐步压缩到 2 导，四分类任务性能出现显著退化，而当任务收束为 C3/C4 双导左右手二分类后，识别能力能够明显恢复。这说明低通道问题的关键不只是结构压缩，更在于让任务目标与采集位置的生理可观测边界重新匹配。

第三，本文在系统层面完成了由离线主线向在线闭环原型的初步收口。基于“预训练初始化 + 轮次级模型更新 + 跨轮次判别反馈”的联动机制，系统已经能够在双导硬件、Qt 上位机和训练桥接脚本之间形成小规模在线闭环。公开数据集上的按时间顺序组织的伪在线结果表明，有限参数重拟合较单纯在线对齐更具有潜力；双导左右手在线闭环试跑则进一步说明系统链路与轮次联动机制已经可以运行，但统计强度仍不足以支持过强的性能主张。总体来看，本文已经初步打通了从标准高密度离线解码到双导端侧原型系统的落地路径。

\section{现存不足}

尽管本文已经完成了从硬件实现、算法设计到系统联动的主体工作，但当前研究仍存在若干明显不足。

首先，在线闭环任务级证据仍然偏少。当前双导左右手在线实验已经进入初步验证阶段，也能够支持预训练初始化、轮次级更新和跨轮次反馈机制的可运行性判断，但无论在被试数量、轮次数量还是统计稳定性方面，仍不足以将其写成强在线性能结论。特别是在更严格的跨轮复盘口径下，预训练初始化的收益尚未稳定优于 baseline，这表明在线闭环场景中的 session-level 漂移和小样本问题仍然突出。

其次，边缘推理节点的系统级验证尚未完全闭环。第 4 章已经给出了模型经剪枝与量化后在 RK NPU 上的推理时间与模型压缩结果，完成了模型侧的端侧部署验证；但当前在线闭环实验的主要执行载体仍为 Linux PC 原型环境，围绕 RK3568 / RKNN 的整链路板端验证、统一时延测量与稳定性统计仍有待补充。

再次，硬件系统的长期运行能力与移动场景能力仍缺乏更完整的量化支撑。虽然当前系统已具备锂电池供电、BLE 连续连接和短时稳定运行的基础，且在现有测试中未见明显持续性异常，但关于长时间运行稳定性、功耗构成、续航时间以及轻度移动场景下的定量性能仍未形成完整实验结果。

最后，当前在线范式仍主要围绕标准化左右手运动想象任务展开。对于低成本双导系统而言，这样的范式有助于先验证系统基础可用性，但尚未充分体现更自然交互、更强个体适配和人机双向协同训练的潜力。因此，本文当前更适合被视为“系统原型与方法主线已经成立”，而不是“应用级脑机接口产品已经完成”。

\section{研究展望}

围绕上述不足，本文后续工作可从以下几个方向展开。

第一，进一步强化在线闭环场景下的人机协同训练机制。相比单纯依赖更大规模模型或更高离线准确率，低通道端侧脑机接口更需要探索如何通过轻量模型、合理交互提示和轮次级在线更新，使被试与模型共同适配。后续可继续研究基于少量参数更新、源数据质量控制和试次质量感知的在线自适应策略，以提升冷启动阶段的稳定性。

第二，将固定、标准化的运动想象范式逐步过渡到更自然的交互任务。当前左右手二分类任务是双导系统的合理切入点，但未来若要走向更完整的实际交互场景，还需要进一步探索如何将运动想象与界面控制、连续意图调节或自然行为反馈结合起来，从而减少纯实验室范式与实际交互之间的割裂。

第三，继续完善端侧原型的完整部署链路。后续研究应将当前 Linux PC 原型中的训练与推理模块逐步迁移到 RK3568 / RKNN 一类嵌入式平台上，完成从模型转换、推理调度到 BLE 数据接入和界面反馈的整链路板端验证，并在统一口径下报告模型时延、系统时延、功耗与稳定性等关键指标。

第四，围绕导联配置与硬件形态开展进一步扩展。本文结果已经表明，当系统仍被限制为双导时，左右手二分类是更稳妥的任务目标；同时，通道组合探索性实验也提示 Cz 可能是未来最值得优先补入的导联。因此，后续可沿着“双导原型 -> 三导增强原型”的路线继续推进，在保持便携性与低成本的前提下进一步改善系统可观测性。

